% !PRINTED 202
\weberpage{202}
%<<CONT
\weq{7}{\begin{gathered}
a_2' = a_0\,(V - W),\\
a_2'' = a_0\,(W - U),\\
a_2''' = a_0\,(U - V).
\end{gathered}}

Kennt man die $a_2', a_2'', a_2'''$, so kann man auch die $U, V, W$
bestimmen durch
\weq{8}{\begin{gathered}
3\,a_0\,U = a_2''' - a_2''\\
3\,a_0\,V = a_2' - a_2'''\\
3\,a_0\,W = a_2'' - a_2',
\end{gathered}}
woraus hervorgeht, dass, wenn von den drei Grössen $a_2', a_2'', a_2'''$
zwei einander gleich sind, eine der Grössen $U, V, W$ verschwindet
, also zwei der Wurzeln der biquadratischen Gleichung
einander gleich sind, und folglich ihre Discriminante verschwindet.

Setzen wir noch
\weq{9}{\begin{gathered}
\alpha\beta + \gamma\delta = u,\\
\alpha\gamma + \delta\beta = v,\\
\alpha\delta + \beta\gamma = w,
\end{gathered}}
so ist
\[
a_0\,(u + v + w) = a_2,
\]
[die Summe der Producte der Wurzeln von $f(x,\,y)$ zu je zweien],
\[
\begin{gathered}
U = v - w,\\
V = w - u,\\
W = u - v,
\end{gathered}
\]
also
\weq{10}{\begin{gathered}
3\,a_0\,u = a_2 - a_2',\\
3\,a_0\,v = a_2 - a_2'',\\
3\,a_0\,w = a_2 - a_2'''.
\end{gathered}}

Es sind also mit den Grössen $a_2', a_2'', a_2'''$ zugleich die $u, v, w$
bekannt. Wenn man aber die Grössen $u, v, w$ kennt, so ist
die Lösung der biquadratischen Gleichung auf Quadratwurzeln
zurückgeführt. Denn es ist
\[
\alpha\beta + \gamma\delta = u, \quad a_0\,\alpha\beta\,.\,\gamma\delta = a_4
\]
und daraus
\[
2\,\alpha\beta = v_1 + \sqrt{\frac{a_0\,u^2 - 4\,a_4}{a_0}}, \quad
2\,\gamma\delta = v_1 - \sqrt{\frac{a_0\,u^2 - 4\,a_4}{a_0}}.
\]
% [resolved 2026-08-06, 5x zoom strip 5/5] v_1 is what is printed, in BOTH halves of this
% display, and the subscript is unambiguous at zoom. Weber introduces v without a subscript
% earlier, so the inconsistency is the author's, not the transcription's -- reproduced as
% printed, per the standing rule on his own oddities.

Da man nun ebenso die Producte
\[
\alpha\gamma, \quad \alpha\delta, \quad \beta\gamma, \quad \beta\delta
\]
%CONT>>
